% !PRINTED 483
\weberpage{483}

gesetzt ist, so dass also die $p, q$ Functionen der Wurzeln $\lambda_r$ sind. Nun wählen wir drei verschiedene Functionen $\eta_0$, die den bisher ausgesprochenen Bedingungen genügen, und bezeichnen sie mit $\eta_0, \eta_0', \eta_0''$.

Die aus diesen drei Functionen abgeleiteten Formen $\psi$ seien $\psi, \psi', \psi''$, und deren Coëfficienten (5) $p_i, q_i$; $p_i', q_i'$; $p_i'', q_i''$. Dann ist nicht nur jede der drei Functionen $\psi, \psi', \psi''$ eine simultane Invariante der Gruppen $P_{168}, G_{168}$, sondern auch ihre Functionaldeterminante (Bd.\,I, \S.\,59):
\weq{6}{\Psi = \frac{1}{S}\begin{vmatrix} \dfrac{\partial\psi}{\partial x_1}, & \dfrac{\partial\psi}{\partial x_2}, & \dfrac{\partial\psi}{\partial x_3}\\[2mm] \dfrac{\partial\psi'}{\partial x_1}, & \dfrac{\partial\psi'}{\partial x_2}, & \dfrac{\partial\psi'}{\partial x_3}\\[2mm] \dfrac{\partial\psi''}{\partial x_1}, & \dfrac{\partial\psi''}{\partial x_2}, & \dfrac{\partial\psi''}{\partial x_3} \end{vmatrix}.}

Die Determinante $\Psi$ ist eine Form 3\textsuperscript{ten} Grades in den Variablen $x$, die nach der Bezeichnungsweise Bd.\,I, \S.\,15, (4) den Ausdruck haben mag:
\weq{7}{\Psi = \sum A_{h,i,k}\, x_h x_i x_k.}

Die Coëfficienten $A_{h,i,k}$ sind dann Functionen der Wurzeln $\lambda_r$, die linear und homogen von jedem der drei Systeme $\eta_r, \eta_r', \eta_r''$ abhängen. Solche Functionen nennt man \gesperrt{trilinear}.

Es bedeute nun $\xi_1, \xi_2, \xi_3$ ein System zu $(x)$ contragredienter Variablen, so dass die biquadratische Form
\weq{8}{f(\xi_1, \xi_2, \xi_3) = \xi_1^3\xi_3 + \xi_2^3\xi_1 + \xi_3^3\xi_2}
eine Contravariante von $G_{168}$ ist. Wenn wir nun in (7)
\[
x_h x_i x_k \quad \text{durch} \quad \frac{1}{6}\frac{\partial^3 f}{\partial\xi_h\partial\xi_i\partial\xi_k}
\]
ersetzen, so erhalten wir eine lineare Form
\weq{9}{L = \frac{1}{6}\sum A_{h,i,k}\,\frac{\partial^3 f}{\partial\xi_h\partial\xi_i\partial\xi_k} = X_1\xi_1 + X_2\xi_2 + X_3\xi_3,}
in der die $X_1, X_2, X_3$ Functionen von $\lambda_r$ sind, und $L$ bleibt nach dem Satze \S.\,129, 3. ungeändert, wenn die Wurzeln $\lambda_r$ durch irgend einer Permutation $\pi$ der Gruppe $P_{168}$ und gleichzeitig die Variablen $(\xi)$ mit den Variablen $(x)$ durch die entsprechende Substitution contragredient transformirt werden. % [sic] \textquotedbl{}irgend einer Permutation\textquotedbl{} as printed
